(a) The assertion is immediate if $n=0$. Otherwise, write $\varphi^*\mathcal O_{\mathbb P^m}(1)\cong\mathcal O_{\mathbb P^n}(d)$ by (6.17). Its generating sections are homogeneous polynomials $f_0,\ldots,f_m$ of degree $d$, so $d\ge0$. If $d=0$, they are constants and $\varphi$ is constant.

Suppose $d>0$. A positive-dimensional geometric fibre would be a projective subset of $\mathbb P^n$ disjoint from $V(f_i)$ for a coordinate nonzero at its image. This contradicts the projective dimension theorem (I, 7.2). Thus every fibre is finite, and the dimension theorem gives $\dim\varphi(\mathbb P^n)=n$. In particular, $m\ge n$.

(b) Let $r+1$ be the dimension of the span of the $f_i$. After a target automorphism, $\varphi$ is given by a basis of this span followed by zero coordinates. The basis elements are independent linear combinations of the degree-$d$ monomials. Hence they define a linear projection from $\mathbb P^N$ to $\mathbb P^r$, whose center misses the $d$-uple image because they have no common zero. Thus $\varphi$ is the $d$-uple embedding, followed by this projection, the linear inclusion $\mathbb P^r\subseteq\mathbb P^m$, and a target automorphism. The integer $d$ is uniquely determined by $\varphi^*\mathcal O(1)$.

When the image spans $\mathbb P^m$, this is the stated projection to $\mathbb P^m$. In general the linear inclusion is necessary; for example, a linear embedding $\mathbb P^1\subseteq\mathbb P^2$ has $d=1$ and $N=1$. Finiteness of the fibres was proved in (a).
